\documentclass[a4paper,fleqn]{cas-dc}

\usepackage[numbers,sort&compress]{natbib}
\graphicspath{{figures/}}
\usepackage{graphicx}
\usepackage{booktabs}
\usepackage{multirow}
\usepackage{amsmath,amssymb}
\usepackage{xcolor}
\usepackage{xurl}
\usepackage{float}

\raggedbottom
\emergencystretch=6em
\tolerance=2000
\hbadness=10000
\hfuzz=130pt
\sloppy

\begin{document}
\let\WriteBookmarks\relax
\def\floatpagepagefraction{1}
\def\textpagefraction{.001}
\hypersetup{hidelinks,hypertexnames=false}
\pdfstringdefDisableCommands{\def\\{ }}

\shorttitle{Weakly Supervised Cryptocurrency Risk-State Learning}
\shortauthors{Lou et al.}

\title[mode=title]{Evaluating Weakly Supervised Cryptocurrency Risk-State Learning:\\When Label Fit, Stability, and Semantics Disagree}

\author[1]{Junwen Lou}
\cormark[1]
\ead{312219020216@home.hpu.edu.cn}
\credit{Conceptualization, Methodology, Software, Formal analysis, Investigation, Writing - original draft}

\author[1]{Fukai Zhang}
\fnmark[1]
\ead{zhangfukai@hpu.edu.cn}
\credit{Supervision, Methodology, Writing - review \& editing}

\author[1]{Shuo Tian}
\fnmark[1]
\ead{312220030817@home.hpu.edu.cn}
\credit{Supervision, Validation, Writing - review \& editing}

\affiliation[1]{organization={Henan Polytechnic University},
            city={Jiaozuo},
            state={Henan},
            postcode={454000},
            country={China}}

\cortext[1]{Corresponding author.}
\fntext[fn1]{Fukai Zhang and Shuo Tian contributed equally and share second authorship.}

\begin{abstract}
Cryptocurrency risk-monitoring and automated trading systems often rely on discrete market-state models, but when supervision is heuristic these models are usually selected by how well they reproduce the heuristic itself. That creates a practical blind spot: label fidelity and deployment usefulness need not coincide. This paper studies how weakly supervised cryptocurrency market-state models should be evaluated when clean ground-truth regimes are unavailable. Using Binance spot 1-hour data, we construct past-only proxy labels from return, volatility, moving-average-gap, and drawdown signals, and compare rule-based, static, post-processed static, recurrent, and temporal-convolutional models. The main contribution is an audit-oriented evaluation protocol that separates proxy-label agreement, temporal stability, and ex-post state semantics, so deployable state representations are not judged by label fit alone. Risk semantics are assessed through ordering tests, nominal nonparametric significance checks, and a block-bootstrap robustness check on BTCUSDT. The experiments reveal systematic ranking reversals. On BTCUSDT, TCN-96x4 improves proxy-label macro-F1 over the default GRU, yet switches more often and is weaker on nominal significant risk layering. A stronger static HistGradientBoosting baseline is selected by weighted classification-plus-stability scoring and TOPSIS, but loses 6/21 nominal significant-risk-layering checks relative to the semantic selector. Across assets and threshold variants, the learned states align more consistently with future risk stratification than with return direction. The central finding is that, in weakly supervised risk-state learning, fit to noisy labels and downstream state usefulness are distinct objectives; model selection that ignores this gap can prefer brittle or weakly interpretable state paths.
\end{abstract}

\begin{keywords}
cryptocurrency market \sep risk state learning \sep weak supervision \sep model audit \sep decision support
\end{keywords}

\begingroup
\makeatletter
\let\cas@origHyWarning\Hy@Warning
\def\Hy@Warning#1{}
\maketitle
\hypersetup{pdfauthor={},pdfsubject={}}
\let\Hy@Warning\cas@origHyWarning
\makeatother
\endgroup

\begin{figure*}[!t]
\centering
\includegraphics[width=\textwidth,trim={0 76mm 0 66mm},clip]{fig1.pdf}
\caption{Overview of the weakly supervised risk-state learning and evaluation workflow. Market windows are converted into rule-based weak labels, candidate learners are evaluated along proxy-label fit, temporal stability, and ex-post risk semantics, and the reported state path is interpreted together with possible ranking reversals across evaluation axes.}
\label{fig:overview}
\end{figure*}

\section{Introduction}
Cryptocurrency risk-monitoring, regime-filtering, and automated trading systems increasingly rely on compact market-state models, but their model-selection logic often has a systematic blind spot: candidate models are typically compared by heuristic-label fit before anyone checks whether the resulting state paths are stable and operationally meaningful. Market-state identification is therefore not only a prediction problem but also a selection problem. In contrast to open-ended state discovery, it fixes a small decision space and asks whether a model can distinguish persistent market conditions in a way that is stable over time and reproducible under modest computing resources. Figure~\ref{fig:overview} summarizes the resulting weak-supervision and evaluation workflow. Regime-switching and hidden-state models provide a long-established language for this problem in macroeconomics and finance \cite{hamilton1989nonstationary,rabiner1989hmm,ang2012regime,mamon2007hmmfinance}.

Recent work has applied this perspective to bull/bear detection, multi-state market classification, and machine-learning-based regime recognition in financial markets \cite{oelschlager2023bullbear,wang2021bearbull,nyberg2013bearbull,koukorinis2025regime}. In cryptocurrency markets, related studies report regime effects in joint return dynamics, volatility, illiquidity, and tail risk \cite{figatalamanca2021crypto,ardia2019bitcoinvol,tan2021bitcoinmsgarch,agakishiev2025forecasting}. The present study uses only Binance spot 1-hour candlestick data and a locally reproducible pipeline, but the scientific question is narrower and more general than a reproducibility exercise: how to learn a persistent market-state representation when the available supervision is heuristic, noisy, and only partially aligned with the latent regimes of interest.

The central difficulty is that hourly market states do not come with clean ground-truth labels. Pure rule systems are easy to reproduce but tend to be brittle. Purely unsupervised state discovery is flexible but often drifts away from the decision question that users actually care about. Weak supervision and noisy-label learning provide a useful middle ground when the available supervisory signal is heuristic, partial, or corrupted rather than exact \cite{ratner2020snorkel,zhang2022programmatic,zhou2018weakly,natarajan2013noisy,frenay2014labelnoise,song2023noisylabels}. In this setting, the target should not be the heuristic label itself, but the latent persistent regime that the heuristic label only imperfectly reveals. We therefore construct coarse proxy labels from past-only trend, volatility, and drawdown statistics, and train temporal models to distill these noisy local cues into a more stable state representation.

This paper makes three contributions. First, it recasts hourly cryptocurrency market-state identification as weakly supervised risk-state learning, so the learning target is a persistent latent state path rather than exact heuristic-label replication. Second, it formalizes a three-dimensional evaluation protocol that separates proxy-label agreement, temporal stability, and ex-post state semantics; within the semantic axis, it further distinguishes raw ordering, nominal significant risk layering, and block-bootstrap robustness. Third, it shows through multi-asset tests, temporal-backbone comparisons, stronger static baselines, deconfounding checks, ablations, uncertainty analysis, overlap controls, and threshold sensitivity that these axes do not induce the same ranking. Post-processed static baselines dominate label fit, a tuned TCN can exceed the default GRU on proxy-label macro-F1, and the nominal significance ranking on BTCUSDT becomes much less decisive once serial dependence is partly respected. The broader methodological point is that fit to noisy labels and downstream risk usefulness are distinct objectives. The protocol turns that distinction into a concrete audit framework for model trustworthiness in financial decision-support settings.

\section{Related Work}
\subsection{Market-State and Regime Identification}
Classical regime-switching research models structural change through latent states and state-dependent dynamics \cite{hamilton1989nonstationary,rabiner1989hmm,ang2012regime,mamon2007hmmfinance}. This literature established the vocabulary that still shapes market-state modelling in finance, especially the emphasis on latent regimes, transition probabilities, persistence, and state-conditioned moments. Building on this foundation, later work studied bull/bear identification and broader multi-state market classification through dynamic binary models, hierarchical hidden Markov models, long-horizon taxonomies, and generative-discriminative classifiers \cite{nyberg2013bearbull,oelschlager2023bullbear,wang2021bearbull,koukorinis2025regime}. These studies are close to our decision problem in spirit, but most focus on equity markets or directional market classification.

Within cryptocurrency markets, related research has documented regime structure in joint return dynamics, Bitcoin volatility, illiquidity, spillovers, and tail risk, often through Markov-switching or regime-dependent risk models \cite{figatalamanca2021crypto,ardia2019bitcoinvol,tan2021bitcoinmsgarch,elkhatib2024illiquid,kayal2024cryptovolnetwork,agakishiev2025forecasting,gao2022caviar,troster2019bitcoinrisk,borri2019tailrisk,koutmos2018liquidity,james2022evolutionary,shakourloo2026globaluncertainty}. More recent machine-learning studies extend this line toward crypto volatility forecasting, driver analysis, and broader reviews of the volatility literature \cite{wang2023cryptovolml,huang2024bitcoinvoldrivers,ahmed2024cryptovolreview}. Most of this work targets volatility prediction, value-at-risk, cross-asset dependence, or directional forecasting rather than learning a stable three-state path directly from spot K-line windows. Our objective is therefore narrower and more operational: to infer whether the resulting states behave as persistent risk strata, not merely as short-horizon return signals.

\subsection{Weak Supervision, Noisy Labels, and State Persistence}
Weak supervision and noisy-label learning provide a natural framework for training from heuristic, partial, or corrupted labels rather than exact annotations \cite{ratner2020snorkel,zhang2022programmatic,zhou2018weakly,hernandez2016taxonomy,schwenker2014partial,natarajan2013noisy,frenay2014labelnoise,song2023noisylabels,tonolini2023wsvae}. This framing matches our setting because the supervisory signal is not hand-labeled ground truth, but a reproducible proxy rule built from past-only return, volatility, moving-average-gap, and drawdown statistics. Recent financial time-series work on adaptive label correction points in the same direction by showing that hand-designed labels can be useful while also making their thresholds, horizons, and correction rules part of the modelling problem \cite{yang2024metalabel}. In that sense, the proxy labels used here should be read as imperfect supervision that may reveal a latent regime only noisily.

A separate but complementary requirement is that the inferred state path should be temporally persistent rather than switching unrealistically often \cite{nystrup2020persistent,bemporad2018jump,livieris2021smoothing,blumke2022structuralhmm}. For market-state models, persistence is not a cosmetic smoothing preference; it is part of what makes the states interpretable and usable as a regime representation. The combination of these two ideas motivates our training target: the model should distill a latent, temporally stable state sequence from heuristic labels rather than merely replicate the label rule bar by bar.

\subsection{Model Evaluation and Temporal Architectures}
Conventional multi-criteria model selection aggregates observed metrics through weighted sums, distance-to-ideal ranking, or related multi-attribute decision rules such as TOPSIS \cite{hwang1981madm}. These methods are useful once the evaluation criteria are already accepted as meaningful objectives. In weakly supervised regime learning, however, the prior question is whether the criteria themselves measure the object of interest. A model can score well on proxy-label agreement and temporal smoothness, yet still produce states whose ex-post economic meaning is weak or unstable. Our evaluation protocol should therefore be read less as a new generic multi-criteria optimizer than as an audit layer that retains conventional classification and stability metrics while adding explicit semantic checks for whether the selected state path behaves like a risk-state taxonomy after the fact. Section~\ref{sec:framework-increment} quantifies this increment against classification-only, weighted classification-plus-stability, and TOPSIS-style classification-plus-stability selectors.

Methodologically, the temporal neural baselines used here are standard rather than architecturally novel. LSTM and GRU encoders remain canonical recurrent models for learning medium- and long-range temporal dependencies \cite{hochreiter1997lstm,cho2014gru}, while temporal convolutional networks provide a non-recurrent alternative based on causal, dilated convolutions and residual blocks that has been shown to compete well on generic sequence modelling tasks \cite{bai2018tcn}. Transformer-style time-series backbones such as PatchTST and iTransformer have become influential in forecasting and representation learning \cite{nie2023patchtst,liu2024itransformer}, but they are left as future baselines because the present study is designed to stress-test the evaluation protocol rather than claim architectural novelty. The experiments therefore focus on whether stronger temporal feature extractors improve only proxy-label fit or also the stability and risk semantics of the inferred state path.

\section{Problem Formulation, Data, and Weak Labels}
\subsection{Problem Formulation}
We formulate hourly cryptocurrency market-state identification as a weakly supervised three-class sequence labeling problem, where the learning target is a temporally persistent latent risk-state path rather than exact replication of heuristic proxy labels. Let $x_t \in \mathbb{R}^{d}$ denote the feature vector extracted from the 1-hour candlestick at time $t$, and let $X_t = [x_{t-w+1}, \dots, x_t]$ be the rolling window of length $w$. The objective is to estimate a three-state regime variable
\begin{equation*}
z_t \in
\begin{cases}
\text{high-risk},\\
\text{transition},\\
\text{low-risk},
\end{cases}
\end{equation*}
with outputs represented as class probabilities
\begin{align*}
&P(z_t=\text{high-risk}),\\
&P(z_t=\text{transition}),\\
&P(z_t=\text{low-risk}).
\end{align*}

The learning target is not a manually labeled ground truth. Instead, we construct a proxy label $\tilde{z}_t$ from past-only return, moving-average gap, volatility, and drawdown signals. The current codebase still emits the historical proxy names bear, neutral, and bull for comparability with the earlier draft, but the present manuscript treats those names as provisional heuristics whose ex-post semantics must be tested rather than assumed. The learning problem is therefore weak supervision \cite{zhou2018weakly,ratner2020snorkel,zhang2022programmatic}: the model should distill a latent, temporally persistent regime path from $\tilde{z}_t$ rather than merely memorize the rule that generated it.

\subsection{Data Source and Split}
The experiments use Binance spot 1-hour candlestick data from 2021-01-01 to 2026-04-20. BTCUSDT and ETHUSDT are the main temporal-model benchmark assets; BNBUSDT, SOLUSDT, and XRPUSDT are added as a multi-asset robustness set for the weak-label, static, post-processed static, and cross-asset-transfer analyses. The default chronological split is:
\begin{itemize}
\item Training: 2021-01-01 to 2023-12-31.
\item Validation: 2024-01-01 to 2024-12-31.
\item Test: 2025-01-01 to 2026-04-20.
\end{itemize}

All splits are chronological. No random reshuffling is used. The test span is reported both as a whole and by the subperiods 2025-H1, 2025-H2, and 2026-YTD (2026-01-01 to 2026-04-20), because a pooled 16-month result can hide whether the learned semantics survive changing market conditions inside the holdout period.

The three BTCUSDT test subperiods also cover materially different market backgrounds. Using the close-to-close path of the underlying spot series, BTCUSDT gains 13.5\% in 2025-H1, then loses 18.4\% in 2025-H2 and 13.6\% in 2026-YTD, while the corresponding max drawdowns remain severe at roughly -30.9\%, -34.8\%, and -35.6\%. This mix of rebound and drawdown phases is one reason why a single pooled directional-return interpretation is unstable even when ex-post risk ordering remains comparatively persistent.

Across the chronological splits, the default proxy labels are imbalanced. For the BTCUSDT/ETHUSDT benchmark assets, neutral occupies 74.7\%--81.3\% of observations, bear 10.6\%--22.0\%, and bull 3.3\%--10.0\% across train, validation, and test segments. The added assets make the imbalance more visible: under the default rule, SOLUSDT and XRPUSDT have sparse bull states in some splits, and under the strict threshold-sensitivity setting their test bull counts fall to 17 and 14 bars, respectively. This imbalance is one reason why macro-F1, balanced accuracy, semantic ordering, significance checks, and explicit check-coverage flags are reported instead of accuracy alone.

\subsection{Feature Set}
Each hourly bar is transformed into 20 hand-crafted features derived from candlestick fields: log returns over 1, 4, and 24 hours; rolling volatility over 24 and 72 hours; high--low range; open--close change; 24-hour z-scores of volume, quote volume, and trade count; exponential moving-average gaps over 12--48 and 24--72 hours; 24-hour rolling skewness and kurtosis; 24-hour up-bar and down-bar ratios; 72-hour maximum drawdown; 24-hour trend strength; taker-buy ratio; and 24-hour volume--price correlation. A 48-step rolling window is then formed for the temporal models. The inclusion of \texttt{max\_drawdown\_72} deserves advance notice because the same statistic also appears in the weak-label rule below. It is therefore an explicit modelling anchor, not a neutral feature, and later deconfounding checks are used to show how much of the recovered risk semantics depend on that anchor.

\subsection{Proxy Label Construction}
The weak labels are constructed from past-only statistics. Historically, the rule generator emits the names bull, bear, and neutral. In the current implementation, a bar is labeled bull when
\begin{align*}
\texttt{log\_return\_24} &> 0.01, \\
\texttt{log\_return\_24} &> \texttt{rolling\_vol\_24}, \\
\texttt{ema\_gap\_24\_72} &> 0.002, \\
\texttt{max\_drawdown\_72} &> -0.03.
\end{align*}
It is labeled bear when
\begin{align*}
\texttt{log\_return\_24} &< -0.01, \\
\texttt{log\_return\_24} &< -\texttt{rolling\_vol\_24}, \\
\texttt{ema\_gap\_24\_72} &< -0.002, \\
\texttt{max\_drawdown\_72} &< -0.06.
\end{align*}
Remaining periods are assigned to neutral. This construction avoids forward-looking leakage in the training target. It also reuses the drawdown statistic on both the feature side and the label side, so the downstream taxonomy is partly anchored to observable downside stress instead of being discovered in a fully free-form way. Later semantic analysis shows why the historical names need to be treated cautiously: empirically, the learned states align more consistently with future risk stratification than with future return direction.

\section{Evaluated Temporal Models and Training Objective}
The evaluated temporal models receive a rolling window $X_t$ and output three-class regime probabilities. The default recurrent variant is a single-layer GRU encoder followed by LayerNorm, dropout, and a linear classification head. We also evaluate a vanilla LSTM under the same windowing and optimization protocol, and a causal dilated TCN with residual temporal convolution blocks. These encoders are treated as comparable backbones inside the same weak-supervision pipeline, not as separate problem formulations.

In addition to cross-entropy on the proxy labels, the temporal models can use two auxiliary regularization terms:
\begin{equation}
\mathcal{L}
=
\mathcal{L}_{\mathrm{ce}}
\,+\,
\lambda_{\mathrm{cont}}\mathcal{L}_{\mathrm{cont}}
\,+\,
\lambda_{\mathrm{bal}}\mathcal{L}_{\mathrm{bal}}.
\end{equation}

The class-balance term penalizes degenerate solutions that collapse into a single dominant class. The continuity term discourages excessive probability changes between adjacent windows:
\begin{equation}
\mathcal{L}_{\mathrm{cont}}
=
\frac{1}{T-1}
\sum_{t=2}^{T}
g_t
\left\|
\begin{aligned}
&\mathrm{softmax}(f(X_t))\\
&-\mathrm{softmax}(f(X_{t-1}))
\end{aligned}
\right\|_1,
\end{equation}
where $g_t$ is a volatility-aware gate,
\begin{equation}
g_t = \frac{1}{1 + \alpha \bar{v}_t},
\end{equation}
and $\bar{v}_t$ is the local rolling volatility summary. Low-volatility periods therefore receive stronger smoothing, while high-volatility periods retain more freedom to change state quickly. In the current implementation, this objective is applied as a chronological batch-local approximation over adjacent windows inside ordered batches rather than as a strict sequence-global penalty. It should therefore be read as a local ordered-window smoothness bias, not as an exact full-path regularizer. The empirical value of these auxiliary terms is examined later by ablation and overlap controls rather than assumed a priori.

\section{Experimental Setup and Evaluation Protocol}
\subsection{Baselines}
The comparison uses a compact but stronger baseline set:
\begin{itemize}
\item Rule-based proxy reference: directly reuse the proxy labels as predictions.
\item Logistic regression: a static supervised baseline on the hand-crafted features.
\item Logistic regression with temporal post-processing: rolling probability smoothing plus a switch-margin rule. This is the key low-cost baseline because it directly tests whether a static classifier with simple smoothing is already sufficient.
\item HistGradientBoosting: a stronger non-linear static baseline on the same hand-crafted features, reported both in raw form and with the same temporal post-processing used for logistic regression.
\item Gaussian HMM: a classical temporal state baseline with three latent states mapped back to the proxy-label space by majority voting on the training split \cite{hamilton1989nonstationary,rabiner1989hmm,oelschlager2023bullbear}.
\item Vanilla LSTM and TCN: standard neural temporal backbones evaluated under the same feature set, chronological split, and loss protocol as the GRU \cite{hochreiter1997lstm,cho2014gru,bai2018tcn}.
\end{itemize}

The rule-based line should be interpreted carefully: because it directly reuses the proxy labels, it is a consistency reference rather than a learned competitor. The neural temporal backbones share the same data pipeline and 48-hour window. Multi-seed checks are used later to determine whether the ranking between post-processed static baselines and recurrent encoders is stable or merely a single-seed accident.

\subsection{Implementation Details}
The temporal models use 48-step windows and the 20-feature input described above. GRU and LSTM use one recurrent layer with hidden dimension 64 and dropout 0.1, followed by LayerNorm, dropout, and a linear classification head. The TCN uses causal convolutions with residual blocks; the reported stress test includes 64-by-3, 128-by-3, and 96-by-4 hidden-dimension/layer configurations. Training uses AdamW with learning rate $10^{-3}$, weight decay $10^{-4}$, batch size 256, and 12 epochs. Unless otherwise stated, the checkpoint with the best validation macro-F1 is retained for testing. For the main GRU, the continuity weight is $\lambda_{\mathrm{cont}}=0.15$, the class-balance weight is $\lambda_{\mathrm{bal}}=0.05$, and the gate strength is $\alpha=3.0$. The temporal post-processing baseline uses a 6-bar rolling mean with a switch margin of 0.10. Main-result tables report the seed-42 run, and separate three-seed summaries are used as robustness checks. A later checkpoint-selection diagnostic compares this default against a semantic-aware alternative, but the main benchmark still uses the label-fit-driven selector so that the comparison with earlier baselines remains fixed.

\subsection{Evaluation Metrics}
The evaluation protocol is organized around three categories of metrics. One category covers classification against the proxy labels: macro-F1, balanced accuracy, bull F1, and bear F1. These measure agreement with the heuristic supervisory signal, not direct accuracy against latent ground-truth regimes. A second category covers temporal stability in the predicted sequences through the total number of regime switches, the average state duration, and the daily switching rate. The third covers ex-post semantics, where we test whether the learned states align more consistently with future return ordering or with future risk ordering. We use the term ranking reversal descriptively for cases where one method is stronger on one evaluation axis but weaker on another, and for that reason the protocol does not collapse these axes into a single utility score.

Formally, for asset $a$ and method $m$, the protocol returns
\[
\Pi(m,a)=\big(C(m,a),T(m,a),S(m,a)\big),
\]
where $C(m,a)$ is the proxy-label agreement vector, $T(m,a)$ is the temporal-stability vector, and
\[
S(m,a)=\big(r^{\mathrm{risk}},r^{\mathrm{return}},g^{\mathrm{nom}},g^{\mathrm{boot}}\big)
\]
collects raw risk-order consistency, raw return-order consistency, nominal significant-risk layering, and, when computed, block-bootstrap robustness. The protocol remains multi-objective: it does not aggregate these quantities into a single score because the experiments repeatedly show that such aggregation would hide ranking reversals. Operationally, the protocol should be read as an audit filter rather than a scalar optimizer. A model that dominates on proxy-label fit but fails semantic checks is still informative about label learnability, but it is not an equally credible deployment choice.

For each default asset--method pair, risk-order consistency is computed over 21 checks: seven risk-oriented metrics (future volatility, absolute return, path loss, path gain, path range, and two downside-hit rates) evaluated over the three test subperiods 2025-H1, 2025-H2, and 2026-YTD. A risk check is counted as matched when the descending state means satisfy bear $>$ neutral $>$ bull. Return-order consistency is computed over six checks using future 24-hour and 72-hour returns with the expected order bull $>$ neutral $>$ bear. To move beyond descriptive ordering, we also run Kruskal--Wallis tests and one-sided Mann--Whitney U tests with Holm correction. A check is counted as a nominally significant risk-layering check when the expected risk order holds, the Kruskal--Wallis $p$-value is below 0.05, and all three adjusted pairwise tests support bear $>$ neutral, neutral $>$ bull, and bear $>$ bull. On the BTCUSDT boundary benchmark, we additionally report a block-bootstrap robustness variant in which the three directional pairwise separations are rechecked with 1000 one-sided block-bootstrap resamples of 24-hour blocks. In threshold-sensitivity experiments, a check is excluded when a state is absent in the corresponding asset--period slice; this mainly affects the strict rule on SOLUSDT and XRPUSDT and is reported through check-coverage flags.

These protocol choices reflect a trade-off. The seven risk metrics cover dispersion, absolute movement, asymmetric downside exposure, broader pathwise variation, and threshold-based loss incidence, which together provide a compact but multi-faceted view of ex-post risk. The three test subperiods balance two competing needs: enough observations per slice for nonparametric testing, and enough temporal resolution to reveal whether semantic orderings survive regime changes inside the 16-month test span. Holm correction is used because each check involves three directional pairwise comparisons; it controls the family-wise error rate while remaining less conservative than Bonferroni in this repeated, relatively low-sample setting. Because the observations remain serially dependent within market periods, the Kruskal--Wallis and Mann--Whitney $p$-values should be read as nominal diagnostics under an independence approximation, not as exact finite-sample evidence. The supplementary block-bootstrap variant does not eliminate all dependence concerns, but it does provide a stricter robustness screen for whether the pairwise risk separations survive coarse serial resampling.

\paragraph{Definition: Semantic regret.}
For a given model-selection rule and candidate pool, semantic regret is the loss in nominal significant-risk-layering rate incurred by the selected model relative to the best-performing model on that semantic axis. It quantifies the hidden deployment cost of favoring proxy-label fit or stability when those criteria point to a model with weaker downstream risk semantics. In practical terms, it measures how much state meaning is lost when model choice is compressed too early into a single score.

\section{Results}
\subsection{BTCUSDT Benchmark and Baseline Boundary}
Table~\ref{tab:btc-compact} establishes the central empirical tension of the paper: the model that best matches the proxy labels is not necessarily the one that produces the most credible state path. These classification scores measure agreement with the heuristic labels rather than direct regime accuracy, so they need to be read alongside switching behavior and the later semantic checks. Raw HistGradientBoosting nearly saturates heuristic-label replication at macro-F1 0.9945, yet it still switches 0.9011 times per day; that profile looks more like near-rule imitation than regime extraction. Once simple temporal post-processing is added, HistGradientBoosting remains the strongest label-fitting learned baseline, with macro-F1 0.8731, a daily switching rate of 0.4084, and an average state duration of 58.47 hours. The post-processed linear baseline is weaker on label fit, but it remains informative because it is cheaper and more transparent. The neural models make the trade-off sharper. Default GRU and vanilla LSTM are smoother than the raw static models, yet both stay below the post-processed static baselines on proxy-label classification. TCN-96x4 reaches macro-F1 0.6817 and moves ahead of the default GRU, but it does so with substantially more switching. Even before semantics enter the discussion, label fit, temporal smoothness, semantic robustness, and model complexity are already pulling in different directions.

The Gaussian HMM is most informative as a failure mode. After the three latent states are mapped back to the proxy-label space by majority voting, all three BTCUSDT hidden states collapse to the neutral label on the test split, with 7006, 2333, and 2062 observations assigned to neutral from the three latent partitions. That is why the reported bull and bear F1 scores are both zero even though the HMM still internally produces three latent-state IDs. In practice, the model is separating something, but not the taxonomy that matters here.

The three-seed checks do not overturn this ordering. GRU reaches macro-F1 $0.6653 \pm 0.0126$ with daily switching rate $0.5291 \pm 0.0088$, whereas LSTM reaches macro-F1 $0.6349 \pm 0.0037$ with daily switching rate $0.5095 \pm 0.0084$. Notably, the gap is not large enough to invite strong claims about backbone superiority in general, but it is stable enough to suggest that the classification-and-stability ranking is not just a single-seed accident. A similar signal appears on the static side in ETHUSDT: raw HistGradientBoosting reaches macro-F1 0.9942, its post-processed version reaches 0.8648, and post-processed logistic regression reaches 0.7391. The stronger non-linear baseline is therefore not merely a BTC-specific curiosity.

\begin{table*}[t]
\caption{Condensed BTCUSDT benchmark. The rule-based proxy itself is omitted because it is a non-learned consistency reference.}
\label{tab:btc-compact}
\centering
\scriptsize
\setlength{\tabcolsep}{4pt}
\begin{tabular}{lccccc}
\toprule
Method & Macro-F1 & Bull F1 & Bear F1 & Switch/day & Avg. Duration \\
\midrule
Logistic regression & 0.7697 & 0.6713 & 0.7351 & 1.2400 & 19.32 \\
Logreg + post-proc. & 0.7435 & 0.6356 & 0.7015 & 0.5474 & 43.68 \\
HistGB & 0.9945 & 0.9973 & 0.9879 & 0.9011 & 26.58 \\
HistGB + post-proc. & 0.8731 & 0.8186 & 0.8395 & 0.4084 & 58.47 \\
Gaussian HMM & 0.2990 & 0.0000 & 0.0000 & 0.0000 & 11401.00 \\
Main GRU & 0.6651 & 0.6251 & 0.4598 & 0.5221 & 45.79 \\
Vanilla LSTM & 0.6314 & 0.5338 & 0.4602 & 0.5179 & 46.16 \\
TCN-96x4 & 0.6817 & 0.6327 & 0.4975 & 0.6779 & 35.30 \\
\bottomrule
\end{tabular}
\end{table*}

The temporal-backbone comparison also shows that the TCN story is not monotonic. The first run, TCN-64x3, improves over LSTM but not over GRU; the wider 128x3 variant then underperforms both; only the deeper 96x4 variant becomes the strongest neural classifier. That matters. It suggests that simply making the TCN wider is not automatically beneficial, whereas a moderate increase in depth helps proxy-label fit under this setup. The trade-off persists, however: all tested TCN variants switch more frequently than the recurrent backbones.

\begin{figure*}[t]
\centering
\includegraphics[width=\textwidth]{btc_temporal_backbone_comparison_extended.pdf}
\caption{BTCUSDT temporal-backbone stress test. TCN-96x4 improves proxy-label macro-F1 over the default GRU, but all tested TCN variants have higher daily switching rates. This result motivates evaluating state paths beyond proxy-label fit alone.}
\label{fig:temporal-backbones}
\end{figure*}

\subsection{Best-TCN Tradeoff: Higher Label Fit, Weaker Risk Semantics}
Among the tested TCN variants, TCN-96x4 gives the highest macro-F1, so we retain it for further comparison with the default GRU. TCN-96x4 improves macro-F1 from 0.6651 to 0.6817, but its daily switching rate rises from 0.5221 to 0.6779. More importantly, the nominal significant risk-layering rate falls from 16/21 checks (76.2\%) to 10/21 checks (47.6\%). Under the present protocol, the stronger TCN reproduces the weak labels more closely, yet produces an ex-post risk-state path that is less stable and less coherent. This dependence on semantic criteria is not a weakness of the analysis; it is exactly the disagreement that the protocol is meant to expose.

\begin{figure*}[t]
\centering
\includegraphics[width=0.74\textwidth]{btc_gru_vs_tcn96_tradeoff.pdf}
\caption{GRU versus the best TCN configuration on BTCUSDT. TCN-96x4 improves proxy-label macro-F1, but the default GRU has lower switching and a higher fraction of nominally significant risk-layering checks.}
\label{fig:gru-tcn-tradeoff}
\end{figure*}

\subsection{Risk-State Semantics Are Stronger Than Return Semantics}
If the analysis stopped with the BTCUSDT benchmark, the practical verdict would look simple: HistGradientBoosting with post-processing would be the strongest static learner, the post-processed logistic model the simplest transparent comparator, and TCN-96x4 the strongest neural classifier. The three-dimensional protocol makes that summary incomplete. Once ex-post semantics are added, the question is no longer just which method best reproduces heuristic labels, but which one produces state paths whose downstream meaning survives changes in period, asset, and significance criterion.

Overall BTCUSDT state semantics show that the three-state taxonomy organizes future risk more consistently than future direction. Across the proxy labels, the most interpretable low-cost baseline (logistic regression with temporal post-processing), and the main GRU, the cleanest separation appears on future 24-hour volatility: the historical bear state is highest risk, bull is lowest risk, and neutral lies between them. The return side is much less stable. In the proxy labels, and at times in the learned models as well, the historical bear state can still be followed by positive average returns.

The directional-return slices show why that story breaks down. Once the BTCUSDT test period is split into 2025-H1, 2025-H2, and 2026-YTD, the ordering of bear, neutral, and bull on future returns changes across subperiods. Figure~\ref{fig:btc-return-semantics} shows the 72-hour return slice, and the full 24-hour plus 72-hour protocol is why all five BTC methods, including the stronger HistGradientBoosting baseline, end up at 0/6 overall return-order consistency checks. The holdout period mixes a rebound phase with two later drawdown phases, and the historical bear label often behaves more like a post-selloff high-risk state than like a persistent negative-return forecast. Future volatility tells a different story: the bear/high-risk state remains consistently above the other two states, while the low-risk state remains lowest. Return semantics vary across subperiods, whereas risk semantics remain more stable.

\begin{figure*}[t]
\centering
\includegraphics[width=0.88\textwidth]{btc_semantics_overall_main.pdf}
\caption{Overall BTCUSDT ex-post state semantics for the proxy labels, the post-processed logistic baseline, and the main GRU. The historical bear/neutral/bull names are retained for comparability, but the observed ordering is more consistent with high-risk/transition/low-risk semantics than with stable directional return forecasting.}
\label{fig:btc-overall-semantics}
\end{figure*}

\begin{figure*}[t]
\centering
\includegraphics[width=\textwidth]{btc_return_semantics_by_period_main.pdf}
\caption{BTCUSDT 72-hour directional-return semantics by subperiod. Across the proxy labels, the post-processed logistic baseline, and the main GRU, the ordering on future 72-hour returns changes across 2025-H1, 2025-H2, and 2026-YTD. The directional bull/bear interpretation is therefore not stable.}
\label{fig:btc-return-semantics}
\end{figure*}

\begin{figure*}[t]
\centering
\includegraphics[width=\textwidth]{btc_risk_semantics_by_period_main.pdf}
\caption{BTCUSDT risk semantics by subperiod. The future-volatility ordering remains much more stable than the return ordering: the bear/high-risk state stays highest, while the bull/low-risk state stays lowest across the three subperiods.}
\label{fig:btc-risk-semantics}
\end{figure*}

The five-asset comparison in Table~\ref{tab:semantic-consistency} makes the same point in a broader setting. The first semantic column reports raw risk-order consistency, the second reports return-order consistency, and the third reports the fraction of 21 risk checks that remain nominally significant under the nonparametric protocol. Notably, risk semantics are much stronger than return semantics across most asset--method pairs. The rankings, however, refuse to settle into a single order: on BTCUSDT, the post-processed logistic baseline has the best raw risk-order consistency, HistGradientBoosting with post-processing dominates proxy-label fit, TCN-96x4 has the best neural proxy-label classification score, and the GRU records more nominally significant risk-layering checks than the other tested temporal backbones. The expanded assets offer both support and a warning. BNBUSDT and SOLUSDT preserve strong risk layering for the lightweight learned baselines, whereas XRPUSDT is materially weaker, especially for direct BTC-trained GRU transfer. The XRP weakness is not only a transfer problem: even the XRP proxy labels themselves reach just 57.1\% raw risk-order consistency and 5/21 nominal significant checks, suggesting that the underlying heuristic taxonomy is already much noisier on that asset. The bootstrap uncertainty intervals reported in the supplementary experiment tables should therefore be read cautiously; they are better viewed as robustness intervals than as high-precision estimates because each rate is based on a small number of period--metric checks.

\begin{table*}[t]
\caption{Three-dimensional semantic evaluation. Risk-order consistency counts how often the expected order bear $>$ neutral $>$ bull appears across 21 risk checks; return-order consistency counts how often bull $>$ neutral $>$ bear appears across six return checks; significant risk checks count how often the expected risk order is also supported by Kruskal--Wallis and Holm-corrected one-sided Mann--Whitney tests under the nominal independence approximation. BTCUSDT block-bootstrap robustness is reported separately in Figure~\ref{fig:btc-bootstrap-robustness}.}
\label{tab:semantic-consistency}
\centering
\scriptsize
\setlength{\tabcolsep}{3pt}
\begin{tabular}{llccc}
\toprule
Asset & Method & \shortstack{Risk-order\\consistency} & \shortstack{Return-order\\consistency} & \shortstack{Sig.\\risk checks} \\
\midrule
BTC & Proxy label & 81.0\% & 0.0\% & 11/21 (52.4\%) \\
BTC & Logistic regression & 90.5\% & 0.0\% & 15/21 (71.4\%) \\
BTC & Logreg + post-proc. & 95.2\% & 0.0\% & 14/21 (66.7\%) \\
BTC & HistGB + post-proc. & 76.2\% & 0.0\% & 10/21 (47.6\%) \\
BTC & Main GRU & 90.5\% & 0.0\% & 16/21 (76.2\%) \\
BTC & Vanilla LSTM & 71.4\% & 0.0\% & 12/21 (57.1\%) \\
BTC & TCN-96x4 & 76.2\% & 0.0\% & 10/21 (47.6\%) \\
\midrule
ETH & Proxy label & 66.7\% & 33.3\% & 9/21 (42.9\%) \\
ETH & Logistic regression & 81.0\% & 33.3\% & 12/21 (57.1\%) \\
ETH & Logreg + post-proc. & 76.2\% & 50.0\% & 13/21 (61.9\%) \\
ETH & Main GRU (ETH-trained) & 71.4\% & 16.7\% & 10/21 (47.6\%) \\
ETH & BTC-trained GRU & 76.2\% & 0.0\% & 9/21 (42.9\%) \\
\midrule
BNB & Proxy label & 66.7\% & 66.7\% & 12/21 (57.1\%) \\
BNB & Logistic regression & 95.2\% & 33.3\% & 15/21 (71.4\%) \\
BNB & Logreg + post-proc. & 95.2\% & 33.3\% & 17/21 (81.0\%) \\
BNB & BTC-trained GRU & 57.1\% & 0.0\% & 7/21 (33.3\%) \\
\midrule
SOL & Proxy label & 66.7\% & 0.0\% & 9/21 (42.9\%) \\
SOL & Logistic regression & 90.5\% & 0.0\% & 14/21 (66.7\%) \\
SOL & Logreg + post-proc. & 90.5\% & 0.0\% & 14/21 (66.7\%) \\
SOL & BTC-trained GRU & 76.2\% & 16.7\% & 13/21 (61.9\%) \\
\midrule
XRP & Proxy label & 57.1\% & 0.0\% & 5/21 (23.8\%) \\
XRP & Logistic regression & 61.9\% & 0.0\% & 6/21 (28.6\%) \\
XRP & Logreg + post-proc. & 71.4\% & 0.0\% & 6/21 (28.6\%) \\
XRP & BTC-trained GRU & 23.8\% & 0.0\% & 0/21 (0.0\%) \\
\bottomrule
\end{tabular}
\end{table*}

The block-bootstrap check gives a more conservative result. Once the three directional pairwise checks are re-evaluated with 24-hour block bootstrap resampling, all methods shrink sharply: the proxy labels fall from 11/21 nominal significant checks to 3/21 bootstrap-robust checks, logreg + post-processing falls from 14/21 to 3/21, the main GRU falls from 16/21 to 3/21, HistGB + post-processing falls from 10/21 to 4/21, and TCN-96x4 falls from 10/21 to 2/21. These results indicate that the nominal counts are useful for ranking models within the protocol, although this should be interpreted cautiously since they are not high-confidence finite-sample evidence once serial dependence is taken more seriously. The ranking among the stronger BTC models also becomes less stable under the stricter screen: the main GRU is no longer uniquely favored. This reinforces the broader point that even the significance layer is better used diagnostically than reduced to a final scalar winner.

\begin{figure*}[t]
\centering
\includegraphics[width=\textwidth]{btc_significance_nominal_vs_block_bootstrap.pdf}
\caption{BTCUSDT nominal versus block-bootstrap significant risk-layering rates. All methods lose a large share of nominal significant checks once 24-hour block bootstrap resampling is used for the directional pairwise separations, and the between-method gaps compress substantially.}
\label{fig:btc-bootstrap-robustness}
\end{figure*}

\subsection{Increment Over Conventional Multi-Criteria Evaluation}
\label{sec:framework-increment}
The preceding results motivate a more explicit comparison with existing multi-criteria evaluation practice. Once different methods lead on different axes, a conventional MCDM pipeline will usually compress that disagreement into a single score and may thereby hide the semantic loss induced by the chosen weighting rule. To quantify that risk, we compare the proposed semantic protocol against four selection rules. One uses classification-only selection by proxy-label macro-F1. Another uses a conventional weighted multi-criteria score over classification and stability, with normalized macro-F1, balanced accuracy, and inverse daily switching rate weighted 0.4, 0.3, and 0.3. A third applies TOPSIS to the same classification-plus-stability criterion set. The last is an oracle-style weighted score that also includes nominal significant-risk layering; it is included to test whether merely adding a semantic scalar is enough, as opposed to treating semantics as a separate diagnostic axis. For this comparison only, the protocol is operationalized as a semantic selector that first maximizes nominal significant-risk layering, then uses lower switching and higher macro-F1 as tie-breakers. This selector is not meant to replace the paper's multi-objective reporting; it is a device for measuring selection disagreement and semantic regret, defined here as the selected method's lost nominal significant-risk-layering rate relative to the protocol selector.

Two candidate pools help clarify where the protocol actually adds something. The five-asset common pool contains logistic regression, logistic regression with post-processing, and the relevant GRU line on each asset. In that pool, classification-only selection disagrees with the semantic protocol on every asset because it always favors the raw logistic model, despite its much higher switching. Once stability is added, most of that discrepancy disappears: weighted scoring and TOPSIS disagree with the protocol on only 20\% of assets, with a mean semantic regret of 0.019 significant-risk rate, or about 0.4 of 21 checks. When the candidate set is simple, the increment over standard classification-plus-stability MCDM is fairly modest.

The BTCUSDT extended-pool comparison summarized in Table~\ref{tab:framework-increment} gives a different result. Once the stronger HistGradientBoosting and TCN candidates enter, classification-only, weighted classification-plus-stability, and TOPSIS all select HistGB with post-processing because it combines the highest macro-F1 (0.8731) with the lowest switching rate (0.4084). Yet under the semantic protocol, that choice loses 0.2857 significant-risk rate relative to the main GRU, which is exactly 6 of 21 risk checks. Even the oracle-style scalar that includes a semantic term does not recover the GRU; it selects the post-processed logistic baseline and still loses 2 of 21 significant-risk checks. In this setting, a method can look reasonable under conventional MCDM and still miss the stronger ex-post risk semantics. This is where the protocol becomes most consequential.

\begin{table*}[t]
\caption{Increment over conventional multi-criteria evaluation. Disagreement is the fraction of assets in the pool where the framework selects a different method from the semantic protocol selector. Semantic regret is the selected method's lost nominal significant-risk-layering rate relative to the protocol selector; 0.2857 corresponds to 6/21 checks. Positive macro-F1 cost means the conventional framework selected a model with higher proxy-label macro-F1 than the protocol selector. Switch cost is the selected method's daily switching rate minus the protocol selector's rate.}
\label{tab:framework-increment}
\centering
\scriptsize
\setlength{\tabcolsep}{4pt}
\begin{tabular}{llcccc}
\toprule
Candidate pool & Framework & Disagreement & Semantic regret & Macro-F1 cost & Switch cost \\
\midrule
Five-asset common & C only & 100\% & 0.038 & 0.044 & 0.898 \\
Five-asset common & C+T weighted & 20\% & 0.019 & 0.016 & 0.005 \\
Five-asset common & C+T TOPSIS & 20\% & 0.019 & 0.016 & 0.005 \\
Five-asset common & C+T+S weighted & 20\% & 0.010 & 0.021 & 0.144 \\
\midrule
BTC extended & C only & 100\% & 0.286 & 0.208 & -0.114 \\
BTC extended & C+T weighted & 100\% & 0.286 & 0.208 & -0.114 \\
BTC extended & C+T TOPSIS & 100\% & 0.286 & 0.208 & -0.114 \\
BTC extended & C+T+S weighted & 100\% & 0.095 & 0.078 & 0.025 \\
\bottomrule
\end{tabular}
\end{table*}

\begin{figure*}[t]
\centering
\includegraphics[width=\textwidth]{framework_increment_semantic_regret.pdf}
\caption{Selection disagreement and semantic regret relative to the semantic protocol selector. Conventional classification-plus-stability MCDM is nearly sufficient in the simple five-asset common pool, but it fails in the BTCUSDT extended pool after stronger label-fitting candidates enter.}
\label{fig:framework-increment}
\end{figure*}

\subsection{Checkpoint Selection as an Audit Problem}
The same tension appears one step earlier, at checkpoint selection time. Up to this point, all temporal models have still used the best validation macro-F1 checkpoint. That is a sensible default for supervised sequence classification, but it is not obviously aligned with the paper's main claim that proxy-label fit is an incomplete deployment criterion. We therefore ran a focused diagnostic on the BTCUSDT GRU. The first selector retained the checkpoint with the highest validation macro-F1. The second used a semantic-aware lexicographic rule on the validation split: it prioritized nominal significant-risk-layering rate, then risk-order rate, then validation macro-F1, and only then lower daily switching rate. A stricter significance-first rule was also tried and immediately selected an immature epoch with much worse test classification, so the reported semantic selector keeps semantic criteria first while still requiring a reasonably mature classifier.

\begin{table*}[t]
\caption{Checkpoint-selection diagnostic for the BTCUSDT GRU. Main window uses the paper's default split (train 2021--2023, validation 2024, test 2025-01 to 2026-04). Rolling window uses train 2021--2022, validation 2023, and test 2024. Significant risk checks report nominal significant-risk layering counts under the same protocol used elsewhere in the paper.}
\label{tab:checkpoint-selection}
\centering
\scriptsize
\setlength{\tabcolsep}{4pt}
\begin{tabular}{llccccc}
\toprule
Window & Selector & Best epoch & Test macro-F1 & Test bal. acc. & Sig. risk checks & Switch/day \\
\midrule
Main & Validation macro-F1 & 12 & 0.6635 & 0.6533 & 14/21 (66.7\%) & 0.5053 \\
Main & Semantic-aware & 9 & 0.6304 & 0.6245 & 16/21 (76.2\%) & 0.5053 \\
\midrule
Rolling 2024 & Validation macro-F1 & 12 & 0.5720 & 0.5340 & 1/14 (7.1\%) & 0.5247 \\
Rolling 2024 & Semantic-aware & 8 & 0.5028 & 0.4983 & 1/7 (14.3\%) & 0.4427 \\
\bottomrule
\end{tabular}
\end{table*}

Table~\ref{tab:checkpoint-selection} shows that semantic-aware checkpointing is neither vacuous nor universally dominant. On the main BTCUSDT window, it changes the selected checkpoint from epoch 12 to epoch 9. That lowers test macro-F1 from 0.6635 to 0.6304, but it raises nominal significant risk layering from 14/21 to 16/21 without increasing daily switching. In other words, the validation split contains enough semantic signal for checkpoint selection to shift the deployment trade-off in the same direction as the paper's broader audit argument.

The rolling-origin window gives a more cautious result. When 2023 is used as validation and 2024 as test, semantic evidence on the validation split is much weaker: the semantic selector never observes a positive nominal significant-risk-layering rate and is forced to rely on weaker ordering signals. In that setting, semantic-aware checkpointing still produces a smoother test path and a slightly higher significant-risk-layering rate, but the classification cost is larger and the gain in ex-post semantics is modest. The practical implication is that semantic-aware checkpointing should itself be audited. It is most useful when the validation window already contains a detectable risk-state structure; when that structure is weak, semantic checkpointing becomes a conservative bias rather than a reliable improvement rule.

\subsection{Threshold Sensitivity and Ablation Checks}
The threshold sweep tests whether the main semantic pattern is merely an artifact of one hand-tuned rule. It does not appear to be. Across all five assets, loosening or tightening the return, EMA-gap, and drawdown thresholds leaves cross-asset mean risk-order rates substantially higher than mean return-order rates. The strict rule also exposes a practical limit of hand-built taxonomies. On SOLUSDT and XRPUSDT, it leaves only 17 and 14 test bull bars, respectively, which reduces some risk checks from 21 to 7 or 14. This result also indicates that pushing the heuristic labels toward apparent precision can easily destroy state coverage.

\begin{table*}[t]
\caption{Threshold-sensitivity summary across five assets. Values are cross-asset means for each variant--method pair. The strict variant has incomplete check coverage for some SOLUSDT/XRPUSDT rows because the bull state becomes nearly absent.}
\label{tab:threshold-sensitivity}
\centering
\scriptsize
\setlength{\tabcolsep}{4pt}
\begin{tabular}{llcccc}
\toprule
Variant & Method & \shortstack{Mean risk-order\\rate} & \shortstack{Mean sig.\\risk rate} & \shortstack{Mean return-order\\rate} & \shortstack{Min risk\\checks} \\
\midrule
Default & Proxy label & 67.6\% & 43.8\% & 20.0\% & 21 \\
Default & Logistic regression & 83.8\% & 59.0\% & 13.3\% & 21 \\
Default & Logreg + post-proc. & 85.7\% & 61.0\% & 16.7\% & 21 \\
\midrule
Loose & Proxy label & 69.5\% & 41.9\% & 23.3\% & 21 \\
Loose & Logistic regression & 75.2\% & 49.5\% & 10.0\% & 21 \\
Loose & Logreg + post-proc. & 75.2\% & 57.1\% & 16.7\% & 21 \\
\midrule
Strict & Proxy label & 68.6\% & 46.7\% & 13.3\% & 7 \\
Strict & Logistic regression & 69.0\% & 54.8\% & 13.3\% & 14 \\
Strict & Logreg + post-proc. & 68.1\% & 50.0\% & 16.7\% & 14 \\
\bottomrule
\end{tabular}
\end{table*}

The drawdown deconfounding check is especially revealing because \texttt{max\_drawdown\_72} sits in both the feature set and the weak-label rule. That creates a legitimate concern: perhaps the recovered risk semantics are partly inherited from a hand-coded downside-stress anchor. Removing the feature weakens both shallow and deep models, which indicates that it carries real signal. But the magnitude of the weakening is uneven. Logistic regression drops from macro-F1 0.7697 to 0.7051 and from 71.4\% to 61.9\% in significant risk layering; the GRU drops from macro-F1 0.6651 to 0.6223 and from 76.2\% to 38.1\%. This pattern suggests that the current GRU leans more heavily on the drawdown feature than the shallow baseline does. That should not be overstated, because the no-drawdown logistic model still retains substantial risk layering. Even so, the drawdown anchor is better treated as an important and transparent design choice than as an incidental feature.

\begin{figure*}[t]
\centering
\includegraphics[width=\textwidth]{btc_drawdown_deconfounding_comparison.pdf}
\caption{BTCUSDT drawdown-feature deconfounding check. Removing \texttt{max\_drawdown\_72} weakens both logistic regression and GRU, but the degradation is larger for the GRU under the current setup.}
\label{fig:drawdown-deconfounding}
\end{figure*}

The complementary rule-side deconfounding check shows the same pattern more clearly. Once the drawdown thresholds are removed from the weak-label rule, classification becomes easier while semantics deteriorate. Logistic regression improves from macro-F1 0.7697 to 0.8778, and the GRU improves from 0.6651 to 0.8194; at the same time, significant risk-layering counts fall from 15/21 to 3/21 for logistic regression and from 16/21 to 7/21 for the GRU. Even the drawdown-free proxy labels themselves retain only 3/21 significant risk-layering checks. This contrast suggests that removing the drawdown anchor does not destroy learnability, but it does weaken the cleanest ex-post risk separation. The present state taxonomy is therefore partly anchored by drawdown-based weak supervision, although this again should be interpreted cautiously rather than as a claim that the entire result is an artifact.

\begin{table*}[t]
\caption{BTCUSDT label-side drawdown-rule deconfounding. Removing the drawdown thresholds from the weak-label rule improves proxy-label fit for both learned models, but it substantially weakens statistically significant risk layering.}
\label{tab:drawdown-rule-deconfounding}
\centering
\scriptsize
\setlength{\tabcolsep}{5pt}
\begin{tabular}{llcc}
\toprule
Weak-label rule & Model & Macro-F1 & Sig. risk checks \\
\midrule
Default & Logistic regression & 0.7697 & 15/21 (71.4\%) \\
Default & Main GRU & 0.6651 & 16/21 (76.2\%) \\
No drawdown filter & Logistic regression & 0.8778 & 3/21 (14.3\%) \\
No drawdown filter & Main GRU & 0.8194 & 7/21 (33.3\%) \\
\bottomrule
\end{tabular}
\end{table*}

The auxiliary-loss controls suggest a more modest interpretation. Removing $\mathcal{L}_{\mathrm{bal}}$ changes the BTCUSDT GRU test macro-F1 from 0.6651 to 0.6628 and the daily switching rate from 0.5221 to 0.5053, while preserving the same significant risk-layering count as the main GRU (16/21). Removing the continuity term leaves macro-F1 almost unchanged at 0.6671 and again preserves 16/21 significant risk-layering checks. Removing the volatility gate changes little as well: macro-F1 0.6655 and 16/21 significant risk-layering checks. The stride-48 no-continuity control is much more damaging: once overlap support is sharply reduced, the BTCUSDT test set shrinks to 238 windows, macro-F1 falls to 0.5794, daily switching rate falls to 0.2996 simply because long overlapping runs are absent, and the significant risk-layering rate collapses to 0/21. This pattern suggests that, under the present setup, the continuity regularizer itself adds only a modest increment on top of highly overlapping windows, whereas overlap reduction substantially weakens the semantic structure recovered by the model. The results therefore do not point to any single auxiliary term as decisive. Instead, they place the temporal objective as one configuration inside a broader protocol whose value is better judged by stability and semantic diagnostics than by a large isolated component effect.

\section{Discussion}
The results do not support a single dominant backbone. Static classifiers with light temporal post-processing remain competitive, and a tuned TCN can outperform the default GRU on proxy-label classification. Better proxy-label fit, however, does not consistently produce smoother or more meaningful state paths. Across Figure~\ref{fig:btc-overall-semantics}, Figure~\ref{fig:btc-return-semantics}, Figure~\ref{fig:btc-risk-semantics}, and Table~\ref{tab:semantic-consistency}, the most persistent pattern is risk ordering rather than directional return ordering. That is the paper's main substantive result: in this setting, weak supervision recovers a more stable risk taxonomy than a directional forecasting taxonomy.

The ranking also depends sharply on the evaluation axis. Post-processed HistGradientBoosting is strongest on label fit, TCN-96x4 is the strongest tested neural classifier on that axis, and the default GRU looks more favorable mainly under the nominal significant-risk criterion on BTCUSDT. Even that advantage becomes less distinct once block-bootstrap checks are added. The main result is therefore not a universal winner but an empirical demonstration that weakly supervised state-learning methods can exchange places once the evaluation target shifts from label imitation to stability and semantics. The methodological contribution is to make that shift inspectable rather than implicit: the protocol exposes when a seemingly strong model owes its appeal mainly to heuristic-label fidelity while giving back too much in stability or semantic reliability.

For knowledge-oriented financial monitoring systems, the weak labels in this study should not be read as the deployed knowledge itself, but as a domain-informed supervisory interface built from transparent market heuristics. The deployable object is the inferred state path and its downstream interpretation in risk monitoring, regime filtering, and decision support. Under that view, the protocol acts as an audit layer between model development and deployment: it helps reject candidates that score well against proxy labels yet encode fragile, unstable, or semantically unreliable market-state representations. The same logic extends beyond cryptocurrency data. Any weakly supervised state-learning system that derives labels from imperfect operational heuristics faces the risk that learning to fit those heuristics will outrun learning to produce a trustworthy state representation.

Within that interpretation, continuity regularization, the balance penalty, and the volatility gate look more like secondary inductive biases than primary drivers of the result. Their estimated effects are modest under the tested settings. More revealing are the comparisons with post-processed logistic regression, HistGradientBoosting, the tuned TCN, and the checkpoint-selection diagnostic, because those comparisons expose how quickly the ranking can reverse once semantics are brought back in. The classical HMM serves a different purpose. It mainly functions as a negative reference, since its mapped three-state partition is often close to neutral.

\subsection{Design Implications}
The results suggest four practical design implications for deployment-oriented financial monitoring systems.
\begin{itemize}
\item Proxy-label fidelity and downstream utility should be audited separately. Strong validation performance on heuristic labels is not a reliable proxy for whether the resulting state path will remain meaningful in use.
\item Stronger temporal extractors are not automatically safer choices. In the present setup, the best TCN improves label fit while increasing semantic regret relative to the main GRU.
\item Transparent modelling anchors matter. The drawdown checks show that explicit and interpretable risk anchors can materially shape whether the learned states remain semantically valid after training.
\item Risk stratification is a more stable target than directional return forecasting. For the weak labels used here, the learned states carry more consistent information about future risk than about future direction, which should guide the design of downstream tasks.
\end{itemize}

\subsection{Limitations}
The present claims are bounded by several limitations. The labels are heuristic rather than ground truth, so proxy-label metrics are not direct measures of regime-identification accuracy. The cross-asset analysis covers only five liquid crypto spot pairs and does not extend to other asset classes or to richer inputs such as order-book, news, and sentiment data. The temporal baseline suite is also more extensive for BTCUSDT than for the added assets. In addition, the drawdown deconfounding checks show that one hand-crafted risk anchor materially affects both the labels and the learned models.

The significance analysis is also incomplete: the cross-asset table relies on nominal nonparametric diagnostics under an independence approximation, and the BTC block-bootstrap analysis shows that those counts can shrink once serial dependence is partly respected. The threshold-sensitivity results likewise show that strict hand-crafted rules can make some states too sparse, especially the bull state for SOLUSDT and XRPUSDT. Checkpoint selection remains only partially resolved. The new diagnostic shows that semantic-aware checkpointing can improve test-time risk layering on the main BTCUSDT window, but it also shows that this benefit depends on the validation window itself containing a detectable semantic signal. The protocol should therefore be read as a compact benchmark and audit framework, not as a definitive market-state standard.

\section{Conclusion}
This paper studies cryptocurrency market states as a weakly supervised risk-state learning problem and proposes a three-axis evaluation protocol for settings without ground-truth regimes. The protocol separates proxy-label agreement, temporal stability, and ex-post state semantics, and on BTCUSDT it further distinguishes nominal significance from block-bootstrap robustness. Across the main experiments, weak supervision recovers persistent state paths whose ex-post interpretation aligns more with future risk than with future return direction.

The central design lesson is that fidelity to noisy labels and usefulness of the learned state path are not interchangeable objectives. Classification-plus-stability selection is often adequate in a simple candidate pool, but it can fail once stronger label-fitting models enter the comparison. In the extended BTCUSDT pool, it loses 6/21 nominal significant risk-layering checks relative to the semantic selector. The checkpoint-selection diagnostic shows that the same trade-off can appear even within one model family: on the main BTCUSDT window, a semantic-aware selector recovers 2 additional significant risk checks without increasing switching, but on an earlier rolling-origin window the semantic signal is too weak for this rule to be reliably advantageous. For developers of crypto risk-monitoring, regime-filtering, and automated trading systems, the main implication is that proxy-label fit alone is an unsafe deployment criterion when the downstream state path must remain stable and interpretable, but semantic checkpointing should itself be conditioned on validation windows that actually exhibit stable risk-state structure. More broadly, the results argue that weakly supervised state-learning systems should be judged by how well they preserve stable and meaningful state representations, not only by how faithfully they reproduce the heuristic labels that trained them.

\clearpage
\printcredits

\section*{Declaration of Competing Interest}
The authors declare that they have no known competing financial interests or personal relationships that could have appeared to influence the work reported in this paper.

\section*{Funding}
This work was supported by the National Natural Science Foundation of China
(Grant No. 62202145) and the Henan Province Science and Technology Research
Project (Grant No. 252102320210).

\section*{Role of the funding source}
The funding sources had no role in study design, data collection, data
analysis, data interpretation, manuscript writing, or the decision to submit
the article for publication.

\section*{Data and code availability}
The raw Binance spot candlestick data used in this study are publicly available
from Binance. Processed datasets, code, experiment scripts, and generated
outputs supporting this study are available upon reasonable request.

\section*{Declaration of generative AI and AI-assisted technologies in the manuscript preparation process}
During the preparation of this work, the authors used generative AI and
AI-assisted technologies to improve language, readability, and manuscript
organization. After using these tools, the authors reviewed and edited the
content as needed and take full responsibility for the content of the published
article.

\clearpage
\bibliographystyle{elsarticle-num}
\bibliography{sn-bibliography}

\end{document}
